\section{Background}
\label{sec:background}

\subsection{Standard ReCom and Its Mixing Time}

The \recom{} Markov chain~\citep{deford2021} operates on the space of
valid $k$-district redistricting plans.
A plan $\pi : V \to [k]$ assigns each census tract $v \in V$ to a district;
$\pi$ is valid if every district is connected and population-balanced
(within tolerance $\epsilon$ of the ideal $P_{\mathrm{ideal}} = \sum_v \mathrm{pop}(v) / k$).

A single \recom{} step:
\begin{enumerate}
  \item Selects a random pair of adjacent districts $(d_i, d_j)$.
  \item Merges them into a super-district $D = d_i \cup d_j$.
  \item Samples a uniformly random spanning tree $T$ of $G[D]$.
  \item Finds a tree edge whose removal splits $T$ into two balanced
        subtrees (i.e., each subtree's population is within $\epsilon$ of
        $P_{\mathrm{ideal}}$).
  \item Accepts the split as the new $(d_i', d_j')$ if one is found;
        otherwise rejects.
\end{enumerate}

The mixing behaviour of \recom{} depends critically on $k$.
For small $k$ (say $k = 8$ or $k = 14$), each step affects a meaningful
fraction of the plan, and the chain explores the plan space efficiently.
For large $k$, the fraction of the plan affected per step is $2/k$, which
is small.
More precisely:

\begin{proposition}[Expected steps to restructure all districts]
\label{prop:recom-mixing}
Under uniform district-pair selection, the expected number of \recom{} steps
before every district has been involved in at least one merge-split is
$\Theta(k \log k)$ by the coupon-collector argument.
\end{proposition}

\begin{proof}
There are $\binom{k}{2}$ pairs of adjacent districts (at most; the actual
number depends on the plan).
Each step selects one pair uniformly at random from the $\approx k$ adjacent
pairs (in a reasonably compact plan, each district has $O(1)$ neighbours).
This is equivalent to the coupon-collector problem with $k$ coupons: the
expected number of trials to collect all $k$ coupons is $k H_k = \Theta(k \log k)$,
where $H_k$ is the $k$-th harmonic number.
\end{proof}

For Texas ($k = 38$), $k \ln k \approx 136$.
Thus roughly 136 steps are needed in expectation just to touch every district
once---and touching every district once is far from sufficient for meaningful
plan-space exploration.
Empirically, autocorrelation in the edge-cut statistic at lag 100 remains
above 0.5 for standard \recom{} on Texas.

\subsection{Coarse-Level Abstraction}

Autry et al.~\citep{autry2021} introduced the key idea that the slow mixing
of fine-level chains can be accelerated by running a coupled chain at a
coarser spatial resolution.

\begin{definition}[Coarsening]
\label{def:coarsening}
Let $G = (V, E)$ be the tract adjacency graph and let $f : V \to C$ be a
surjective map from tracts to counties (the coarsening map).
The \emph{county adjacency graph} $G_C = (C, E_C)$ has vertex set $C$
and edge set
\[
  E_C = \{(c_1, c_2) \in C \times C : c_1 \neq c_2
            \text{ and } \exists\, (v_1, v_2) \in E
            \text{ with } f(v_1) = c_1,\, f(v_2) = c_2\}.
\]
That is, two counties are adjacent in $G_C$ if and only if at least one
tract in the first county is adjacent to at least one tract in the second
county in the fine graph $G$.
\end{definition}

The coarsening $f$ induces a \emph{coarse plan}: given a fine plan $\pi : V \to [k]$,
the coarse plan assigns each county $c$ to the district that contains the
majority of $c$'s population, or more precisely, the county plan
$\pi_C : C \to [k]$ is defined by a projection rule (see
Section~\ref{sec:algorithm}).

A coarse \recom{} step operates on $G_C$ with county populations
$\mathrm{pop}(c) = \sum_{v \in f^{-1}(c)} \mathrm{pop}(v)$.
The result is a coarse plan $\pi_C' : C \to [k]$, which is then
\emph{projected} back to a fine plan $\pi' : V \to [k]$ by assigning each
tract $v$ to the district of its county: $\pi'(v) = \pi_C'(f(v))$.

\subsection{Why Coarse Moves Accelerate Mixing}

A coarse \recom{} step is a macro-move that simultaneously reassigns all
tracts in the merged county-pair.
If the two merged counties together contain $m$ tracts, the coarse step
achieves in one move what would require $O(m)$ fine steps.
Since Texas has $\approx 5200$ census tracts spread across $\approx 100$
active county regions, a single coarse step can in principle restructure up
to $52$ tracts per county (average).

More formally, the mixing time of the combined chain is shorter than the
fine chain alone because coarse moves provide \emph{shortcut paths} in
the plan-space graph:

\begin{proposition}[Shortcut paths via coarse moves]
\label{prop:shortcut}
Let $\pi_0$ and $\pi^*$ be two valid fine plans that differ in $\ell$ tracts,
all within a single pair of counties.
The fine \recom{} chain requires $\Omega(\ell / (n/k))$ steps in expectation
to restructure this region (the probability of selecting the relevant
district pair at each step is $\approx 2 / k$).
A single coarse move restructures the entire county pair in one step,
provided the resulting coarse plan is balanced and the rebalance subroutine
succeeds.
\end{proposition}

The coarse chain thus provides a mechanism for large-scale restructuring
that the fine chain lacks.
The interleaving parameter $\alpha$ controls the balance between
large-scale exploration (coarse moves) and local precision (fine moves).
